% GENERATED by tools/gen.py from the unit's vocabulary -- do not edit.
Ćwiczenia tego bloku są \emph{generowane} z listy słów tygodnia i sprawdzają pamięć słownictwa, nie gramatykę. Jeśli odpowiedź jest synonimem, a nie słowem z listy, sprawdź w tabeli powyżej.
\begin{cwiczeniadalej}
\exitem{Jak to jest po niemiecku: tak że?}{\de{sodass}}
\exitem{Jak to jest po niemiecku: chyba że?}{\de{es sei denn}}
\exitem{Jak to jest po niemiecku: ani … ani?}{\de{weder … noch}}
\exitem{Jak to jest po niemiecku: zarówno … jak i?}{\de{sowohl … als auch}}
\exitem{Jak to jest po niemiecku: bez tego, że?}{\de{ohne dass}}
\exitem{Jak to jest po niemiecku: bez?}{\de{ohne zu}}
\exitem{Jak to jest po niemiecku: zatem, wedle tego?}{\de{demnach}}
\exitem{Jak to jest po niemiecku: natomiast?}{\de{dagegen}}
\exitem{Jak to jest po niemiecku: na przykład?}{\de{beispielsweise}}
\exitem{Jak to jest po niemiecku: ogółem?}{\de{insgesamt}}
\exitem{Co znaczy \de{indem}?}{przez to że}
\exitem{Co znaczy \de{falls}?}{jeśli (w razie gdyby)}
\exitem{Co znaczy \de{es sei denn}?}{chyba że}
\exitem{Co znaczy \de{sowohl … als auch}?}{zarówno … jak i}
\exitem{Co znaczy \de{anstatt dass}?}{zamiast tego, że}
\exitem{Co znaczy \de{ohne dass}?}{bez tego, że}
\exitem{Co znaczy \de{dennoch}?}{mimo to}
\exitem{Co znaczy \de{folglich}?}{w konsekwencji}
\exitem{Co znaczy \de{demnach}?}{zatem, wedle tego}
\exitem{Co znaczy \de{hingegen}?}{natomiast}
\end{cwiczeniadalej}
